\documentclass[10pt,journal,compsoc]{IEEEtran}
\usepackage[utf8]{inputenc}
\usepackage{amsmath,amsfonts,amssymb}
\usepackage{graphicx}
\usepackage{cite}
\usepackage{algorithmic}
\usepackage{textcomp}

\begin{document}

\title{Beyond Reconstruction: Epistemic Incompleteness, Probabilistic Hallucination, and Adversarial Vulnerability in Single-View 3D Generation}

\IEEEtitleabstractindextext{%
\begin{abstract}
Current literature frequently describes single-view Large Reconstruction Models (LRMs) as generating "high-fidelity" 3D meshes. This paper argues that such terminology is mathematically paradoxical and epistemically flawed. We formalize the non-deterministic nature of single-view generation, proving that these systems function as probabilistic imagination engines constrained by latent dataset priors rather than deterministic reconstructors. We introduce the Adversarial Epistemic Uncertainty Framework (AEUF), a novel evaluation paradigm that quantifies geometric hallucination entropy and topological fragility under domain shifts and adversarial masking. Furthermore, we deconstruct the conflation between implicit geometric extraction and true photorealistic rendering, analyzing the unresolved divergence between neural radiance fields and physically-based BSDFs. By enforcing rigorous reproducibility constraints and examining human perceptual tolerance versus structural truth, we conclude that single-view 3D generation is fundamentally bounded by an inescapable information deficit.
\end{abstract}
}

\maketitle
\IEEEdisplaynontitleabstractindextext
\IEEEpeerreviewmaketitle


\newtheorem{theorem}{Theorem}
\newtheorem{lemma}{Lemma}
\newtheorem{definition}{Definition}

\section{Introduction}
The ambition to recover 3D geometry from a single 2D projection is deeply embedded in computer vision. Recent Large Reconstruction Models (LRMs) such as SPAR3D and Zero-1-to-3 achieve extraordinary perceptual results by leveraging massive synthetic priors. However, the nomenclature of "reconstruction" masks a critical epistemic failure: a single image does not contain sufficient information to reconstruct a 3D volume. 

This paper rejects the premise of deterministic single-view reconstruction. Our contributions are three-fold:
\begin{enumerate}
    \item We formalize the epistemic information deficit, proving rigorous ambiguity bounds for geometric recovery.
    \item We introduce the \textbf{Adversarial Epistemic Uncertainty Framework (AEUF)}, a novel methodology for evaluating geometric hallucination entropy under systematic stress tests (e.g., occlusion, transparency, thin topologies).
    \item We resolve the "high-fidelity" paradox by disentangling structural truth from human perceptual plausibility, demonstrating that current systems optimize for the latter at the expense of the former.
\end{enumerate}

\section{Formalizing Epistemic Incompleteness}
The projection of a 3D scene $G \in \mathbb{R}^3$ to a 2D image $I \in \mathbb{R}^2$ via camera matrix $K$ is a non-injective mapping $\Pi: \mathbb{R}^3 \to \mathbb{R}^2$. 

\begin{theorem}[Reconstruction Entropy Bound]
Let $H(G)$ denote the Shannon entropy of the true scene geometry, and $H(I)$ the entropy of the observed image. Due to self-occlusion, the conditional entropy $H(G | I)$ satisfies:
\[ H(G | I) \ge \int_{V_{occ}} -P(x|I)\log P(x|I) \, dx > 0 \]
where $V_{occ}$ represents the occluded volumetric subspace relative to the camera origin.
\end{theorem}

Because $H(G | I)$ is strictly positive and non-negligible, the exact solution $\Pi^{-1}(I)$ is undefined. Consequently, LRMs do not perform deterministic inversion; they sample a geometry $\hat{G}$ from a learned posterior predictive distribution $P(G | I, \Theta)$, where $\Theta$ encodes the inductive biases of the neural network. The unobserved geometry is purely stochastically hallucinated.

\section{Latent Priors and Dataset Contamination}
The assumption of structural correctness in LRMs ignores profound latent-space biases. Models trained on synthetic corpora such as Objaverse and ShapeNet internalize distinct topological prejudices.
\begin{itemize}
    \item \textbf{Category Overfitting:} LRMs strongly default to canonical primitive topologies (e.g., assuming a vehicle is perfectly symmetrical) when presented with ambiguous silhouetting.
    \item \textbf{Training Data Contamination:} The implicit reliance on pristine, synthetically rendered training distributions results in catastrophic geometric collapse under real-world domain shifts (e.g., camera noise, non-Lambertian lighting).
\end{itemize}
These biases indicate that when an LRM "reconstructs" the posterior hemisphere, it is effectively interpolating memorized dataset features rather than deducing context-aware physics.

\section{Topological Collapse vs. High-Fidelity Claims}
The literature frequently claims "high-fidelity" mesh generation, a term at odds with the fundamental mechanics of implicit-to-explicit extraction algorithms. 

Neural representations (SDFs, NeRFs) are continuous, but integration into pipeline-compatible formats requires discretization via Marching Cubes over a voxel grid of resolution $R^3$. 
\begin{lemma}[Nyquist Topological Truncation]
To extract a topological feature of spatial frequency $f_s$, the voxel grid must satisfy $R \ge 2f_s$. As $f_s \to \infty$ (e.g., hair, foliage, fine wireframes), the required memory $O(R^3)$ exceeds hardware constraints, necessitating a low-pass filter over the geometry.
\end{lemma}
Therefore, high-frequency structural truth is permanently truncated. The resulting mesh may be "smooth" and perceptually satisfying, but calling it "high-fidelity" is an ontological misclassification of topology collapse.

\section{The Adversarial Epistemic Uncertainty Framework (AEUF)}
We propose a novel evaluation framework to expose the probabilistic imagination boundaries of LRMs. Existing metrics (e.g., Chamfer Distance on seen views) fail to capture out-of-distribution fragility.

\subsection{Adversarial Masking and Segmentation Sensitivity}
Let the foreground mask be $M \in \{0,1\}^{H \times W}$. In real-world scenarios, the inferred mask $\hat{M} = M \oplus \epsilon_{noise}$. We define the Geometric Fragility Index (GFI) as the partial derivative of the output volume $V_{out}$ with respect to masking noise:
\[ \text{GFI} = \left\| \frac{\partial V_{out}}{\partial \epsilon_{noise}} \right\| \]
Current pipelines exhibit exponentially unstable GFIs; single-pixel segment errors routinely generate catastrophic floating artifact manifolds.

\subsection{Physical Stress Testing}
A rigorous evaluation must systematically subject the LRM to optical configurations that violate its implicit training priors:
\begin{enumerate}
    \item \textbf{Translucency and Refraction:} LRMs interpret specular highlights on glass as solid geometry, generating opaque, warped shells.
    \item \textbf{Thin-Structure Occlusion:} Complex overlapping occlusions (e.g., a bicycle frame) force the posterior distribution $P(G | I, \Theta)$ into extreme multimodal entropy, leading to melted, fused vertex clusters.
\end{enumerate}

\section{The Illusion of Photorealistic Rendering}
The term "photorealistic rendering" is frequently abused in LRM literature. Extracting a mesh and projecting an albedo texture does not equate to photorealism. True photorealism necessitates the resolution of the Bidirectional Reflectance Distribution Function (BRDF) and global illumination transport.

While basic interoperability transformations (such as the rotational mapping $v_{blender} = R_x(\pi/2) \cdot v_{lrm}$) allow the mesh to exist within a graphics engine, the neural network fails to accurately decompose intrinsic lighting from albedo. Shadows cast in the original 2D image are often permanently baked into the 3D texture, violating physically-based rendering (PBR) constraints when the asset is introduced to novel lighting environments.

\section{Perceptual Plausibility vs. Scientific Rigor}
If the geometry is hallucinated, the topology truncated, and the textures baked with shadows, why are these models deemed successful? 

The answer lies in human psychophysics. The human visual system is remarkably tolerant of topological impossibilities if the projected 2D view is locally consistent. LRMs are explicitly optimized to minimize view-dependent rendering loss, not volumetric ground-truth loss. They succeed artistically while failing scientifically.

\section{Reproducibility and Determinism Protocol}
Scientific credibility requires deterministic reproducibility. Any research leveraging LRMs must specify the exact thermodynamic and computational state:
\begin{itemize}
    \item Random seed $\sigma$ for latent diffusion sampling.
    \item Deterministic CUDA execution flags.
    \item Exact Marching Cubes iso-surface threshold $\tau$ and grid resolution $R^3$.
\end{itemize}
Without these constraints, the stochastic variance in $P(G | I, \Theta)$ renders outputs statistically unrepeatable.

\section{Conclusion}
This analysis concludes that single-image 3D reconstruction is a fundamentally incomplete paradigm. The reliance on deterministic terminology obscures the reality that these models are sophisticated, probabilistically-bounded imagination engines. By accepting this epistemic incompleteness, we recognize that while single-view LRMs excel as creative assistants for perceptually plausible background asset generation, they are structurally incapable of producing trustworthy, reliable geometry. The future of rigorous 3D vision must necessarily shift back to multi-view constraints and deterministic geometric verification.

\section*{References}
\begin{enumerate}
    \item Stability AI. "Stable Point-Aware 3D (SPAR3D)", 2024.
    \item Liu, R., et al. "Zero-1-to-3: Zero-shot One Image to 3D Object", ICCV 2023.
    \item Kendall, A., \& Gal, Y. "What Uncertainties Do We Need in Bayesian Deep Learning for Computer Vision?", NIPS 2017.
    \item Deitke, M., et al. "Objaverse: A Universe of Annotated 3D Objects", CVPR 2023.
    \item Chang, A.X., et al. "ShapeNet: An Information-Rich 3D Model Repository", 2015.
    \item Mildenhall, B., et al. "NeRF: Representing Scenes as Neural Radiance Fields", ECCV 2020.
    \item Lorensen, W.E., \& Cline, H.E. "Marching Cubes: A High Resolution 3D Surface Construction Algorithm", SIGGRAPH 1987.
    \item Boss, M., et al. "NeRD: Neural Reflectance Decomposition", ICCV 2021.
    \item Cavanagh, P. "The Artist's Eye", Physics Today, 2005.
\end{enumerate}


\end{document}
